\subsection{Herramientas y procedimiento de verificación}
\label{sec:herramientas}

\begin{table}[H]
\centering\footnotesize
\caption{Herramientas del proyecto.}
\label{tab:herramientas}
\rowcolors{2}{SoftGray}{white}
\begin{tabularx}{\textwidth}{@{}>{\bfseries\raggedright\arraybackslash}p{4.4cm} >{\raggedright\arraybackslash}X@{}}
\toprule
\rowcolor{TableHead} Herramienta & Para qué se usó \\
\midrule
Python 3.12 & Simulación de la misión, modelo de referencia del núcleo y figuras \\
Icarus Verilog 12~\cite{icarus} & Simulación de todos los bancos de prueba \\
Active-HDL 17, edición estudiantil & Proyecto con todos los módulos y macros para ver las formas de onda \\
Quartus Prime Lite 24.1 & Síntesis, colocación, ruteo y análisis de tiempos para la MAX~10 \sig{10M50DAF484C7G} \\
DE10-Lite~\cite{terasic}, tres tarjetas & Pruebas en hardware: una, dos y tres tarjetas unidas por UART \\
Sonda de Quartus por JTAG~\cite{intelissp} & Leer los registros de las tarjetas y fijar sus entradas desde la computadora \\
LibreLane 3.0.17 con la PDK \sig{sky130A}~\cite{librelane,skywater} & Layout del núcleo, de RTL a GDS, en una máquina virtual con Ubuntu 24.04 \\
\bottomrule
\end{tabularx}
\end{table}

Cada bloque se verificó solo antes de conectarlo, en este orden:
\begin{enumerate}
  \item \textbf{Modelo de referencia.} Antes de escribir el Verilog del núcleo se programó en Python un
        modelo que hace las mismas cuentas con enteros, registro por registro. Ese modelo genera
        los cuadros de prueba y la respuesta esperada.
  \item \textbf{Bancos autoverificables.} Cada banco compara contra un valor esperado y termina con una
        línea de resultado; ninguno se revisa a ojo. Los bloques numéricos se comparan contra el
        valor exacto en aritmética real, y la distancia, contra todos los puntos del mapa.
  \item \textbf{Integración.} El banco \sig{tb\_tri} cablea las tres FPGA como en el laboratorio y recorre
        nueve escenarios; lee los displays y los LED como los vería una persona.
  \item \textbf{Síntesis y tiempos.} Las cuatro variantes de la red se compilan con el reloj de 50~MHz
        y se exige holgura positiva en el peor caso y cero advertencias críticas. Después de cada
        cambio se comparan los recursos con la síntesis anterior: una diferencia inesperada delata
        lógica eliminada.
  \item \textbf{Tarjetas.} Los mismos escenarios se repiten en las DE10-Lite fijando las entradas y
        leyendo los registros por la sonda de JTAG, sin instrumentos externos.
  \item \textbf{La terminal sin tarjetas.} Para probar el programa de la computadora cuando no hay
        tarjetas, sus funciones de JTAG se sustituyen por archivos y del otro lado se simulan las
        tres FPGA con sus módulos superiores: la orden que escribe el operador, los bits que llegan
        a la FPGA, lo que el circuito hace y lo que la terminal muestra se comprueban de corrido.
\end{enumerate}

\subsubsection{Flujo del layout}
\label{sec:flujo-layout}

El núcleo de la mochila, que es lo que iría en un circuito integrado, se pasó por el flujo
abierto completo de RTL a GDS con la PDK \sig{sky130A} de SkyWater~\cite{skywater} y la
biblioteca de celdas \sig{sky130\_fd\_sc\_hd}. La herramienta fue LibreLane 3.0.17~\cite{librelane},
el sucesor de OpenLane~\cite{shalan2020}: usa los mismos programas (Yosys para la síntesis,
OpenROAD para el plano, la colocación, el árbol de reloj y el ruteo, y Magic, Netgen y KLayout
para la verificación física) y el repositorio de OpenLane remite a él para diseños nuevos. La
entrada son los cinco archivos Verilog del núcleo, sin modificar.

La configuración tiene tres decisiones. El reloj es de 20~ns, el de 50~MHz de la tarjeta,
porque el generador de estímulo cuenta sus pulsos con esa frecuencia. La ocupación del núcleo
se pidió en 40~\% y la densidad de colocación en 55~\%, para rutear sin congestión. Y se activó
la reparación de transiciones lentas después del ruteo global, con márgenes de 50~\%: con los
valores por omisión la corrida pasaba DRC y LVS pero dejaba 140 transiciones de más de 0.75~ns
en la esquina lenta, causadas por los búferes débiles que agrega la reparación de \emph{hold};
con ese ajuste quedaron en cero. Los resultados están en la sección~\ref{sec:layout}.
